Software and modern web services depend increasingly on cloud infrastructure, third-party APIs, and geographically distributed execution. This has improved reliability and velocity but has degraded observability of resource impact: users cannot see the environmental consequences of individual their actions, and providers often cannot attribute energy and emissions inside virtualised stacks as cloud providers do not disclose reliable metrics.

A recurring idea in sustainable systems is to bring resource accounting closer to the interface where decisions are made: rather than computing annualised estimates, disclose per-unit costs in real time. This is the foundation of the Software Carbon Intensity (SCI) value which accounts per unit of work $R$\cite{gsf-sci}.  In web systems the natural unit is the HTTP request. HTTP already functions as an interoperability membrane across organisations, languages, and infrastructures; it is therefore a plausible focus for environmental disclosure. However, real-time energy measurement of arbitrary cloud workloads is difficult without dedicated hardware access and careful experimental controls, and request handling typically fans out to caches, databases, queues, and external services. While there has been ample work on the idea of adding energy/ carbon values to HTTP requests\cite{martin-http-carbon-emissions-scope-2-00,httpwg-admin-issue-52-carbon-emissions,ietf-http-wg-mailinglist-carbon-header-thread} no workable solution has been presented to date, that we are aware of.

This paper makes four contributions:
(1) a design for per-request impact disclosure via HTTP response headers grounded in established HTTP extensibility mechanisms \cite{rfc9110,rfc8941} and naming guidance \cite{rfc6648};
(2) a measurement-to-model method that converts controlled endpoint benchmarks into compact runtime-evaluable energy models;
(3) an implementation of an NGINX module that reads an endpoint model registry and emits per-request energy and \(\mathrm{CO_2e}\) metadata (operational and embodied components);
(4) an evaluation including overhead measurement


\section{Introduction}



\paragraph{Carbon accounting for software and systems}
A central challenge in green software is aligning engineering metrics with carbon accounting. The Software Carbon Intensity (SCI) specification proposes a rate-based framing: operational emissions (energy times grid intensity) plus embodied emissions, normalised by a functional unit \cite{gsf-sci}. While SCI does not mandate a transport mechanism, its structure strongly motivates per-request disclosure for HTTP APIs, where ``one request'' is a natural functional unit.

Life-cycle assessment standards (e.g., ISO 14040/14044) provide methodology for embodied impacts in manufactured products and infrastructure \cite{iso14040,iso14044}. The GHG Protocol Product Standard similarly frames life-cycle greenhouse gas accounting and emphasises boundary clarity \cite{ghgproduct}. For web services, embodied emissions are often hard to attribute precisely; nevertheless, excluding them can mislead, especially for long-lived infrastructure. T

\paragraph{Measuring and modelling software energy}
Direct energy metering of software has a long research lineage. For example JouleMeter estimates power consumption of software components using resource utilisation models calibrated to hardware \cite{kansal2010joulemeter}. Other measurement and profiling approaches combine hardware counters, external meters, and controlled benchmarking to estimate energy per operation. A consistent theme is that production measurement is challenging: noise, multi-tenancy, and hardware abstraction erode attribution quality.

\paragraph{AI inference and request-scaled impact}
Much early work on ML sustainability focused on training cost, demonstrating that improved accuracy can imply substantial additional compute and emissions \cite{strubell2019energy}. Subsequent work highlights that inference can dominate when request volume is high, and that inference behaviour depends on model architecture and workload characteristics \cite{patterson2021carbon,hoffmannmajuntke2024hotcarbon}. This paper therefore treats AI endpoints (\texttt{/ai}) not as special cases but as first-class HTTP endpoints with parameterisable energy models.

\paragraph{Network energy \& emissions accounting}
End-to-end emissions of online services include network transfer. Estimating network energy per byte is difficult to achieve; a widely cited synthesis is Aslan et al.'s meta-analysis on the electricity intensity of internet data transmission, which documents variability across studies and methodological choices \cite{aslan2018internet}. In parallel, the IETF has produced standards for energy management and monitoring in network equipment (e.g., EMAN framework), which address device-level energy management rather than request-level disclosure \cite{rfc7326}. Our work is complementary: we do not attempt to standardise network-device power management; instead, we standardise a disclosure interface at the HTTP layer, which can incorporate network estimates if available but remains functional without them.

\paragraph{HTTP extensibility and telemetry headers}
HTTP explicitly supports extensibility through header fields \cite{rfc9110}. The IETF has standardised structured representations for header values (Structured Field Values) to improve interoperability of non-trivial metadata \cite{rfc8941}. At the same time, naming conventions matter: the historic \texttt{X-} prefix for experimental headers is deprecated because it impedes later standardisation and causes collisions \cite{rfc6648}. Existing telemetry-style headers (e.g., W3C Server-Timing) demonstrate the practical value of surfacing server-side metrics in responses, while raising familiar caveats about privacy, caching, and interpretation \cite{w3c-servertiming}.

\paragraph{API user (consumer) perspective.} Many organisations now operate internal or regulatory carbon accounting processes. Yet API consumption is typically billed in currency, not in \(\mathrm{CO_2e}\) (or even energy). When an application calls \texttt{/login} or \texttt{/ai}, the client can measure cost, latency but not impact. If impact were disclosed directly in an API response, consumers could (i) budget or cap footprint, (ii) optimise client behaviour (batching, caching, prompt reduction), and (iii) compare providers for functionally equivalent services. This aligns with the logic of per-functional-unit accounting emphasised in software carbon intensity frameworks \cite{gsf-sci}.

\paragraph{Service provider perspective.} Providers face the complementary problem: request-level attribution is hard in multi-tenant clouds. Virtual machines, containers, and serverless platforms abstract away low-level power telemetry; energy data may exist only at coarse rack or cluster granularity, if at all. Independently, providers are incentivised to optimise efficiency and to demonstrate progress, but they lack a \emph{drop-in} mechanism to communicate request-level impact to users. A server-side disclosure mechanism that can run at the HTTP layer can also be used internally (e.g., via access logs) to find hotspots and regressions without deploying power meters to production.


----


\subsection{Measurement setup}
Measurements were conducted on a machine equipped with an Intel Core i5-9600K CPU and an NVIDIA GeForce GTX 1080 GPU. Energy was recorded at three levels simultaneously: the full system power supply (MCP39F511N), CPU and DRAM (Intel RAPL via MSR registers), and GPU (NVIDIA NVML). The LLM used was gemma3:1b, served locally via ollama with GPU passthrough. All runs were orchestrated with the Green Metrics Tool \cite{gmt}, which segments execution into named phases that map directly to benchmark steps; per-request energy is derived by dividing the total phase energy by the request count executed in that phase.

\subsection{Per-request energy results}

\begin{table}[h]
\centering
\caption{Per-request system energy (PSU) measured at the power supply unit.}
\label{tab:energy}
\begin{tabular}{lrr}
\toprule
\textbf{Endpoint / condition} & \textbf{Requests} & \textbf{Energy (mWh/req)} \\
\midrule
\texttt{/logout}                  & 100 & 0.129 \\
\texttt{/getToDos} (100 items)    & 100 & 0.196 \\
\texttt{/getToDos} (1\,000 items) & 100 & 0.250 \\
\texttt{/getToDos} (10\,000 items)& 100 & 0.310 \\
\texttt{/getToDos} (100\,000 items)&100 & 1.300 \\
\texttt{/createToDo} (text, 100 chars)   & 100 & 0.412 \\
\texttt{/createToDo} (text, 1\,000 chars)& 100 & 0.471 \\
\texttt{/createToDo} (text, 10\,000 chars)&100 & 0.381 \\
\texttt{/createToDo} (text, 100\,000 chars)&100 & 0.470 \\
\texttt{/done}                    & 100 & 0.393 \\
\texttt{/createToDo} (file, 1\,KB)  & 100 & 0.374 \\
\texttt{/createToDo} (file, 10\,KB) & 100 & 0.360 \\
\texttt{/createToDo} (file, 100\,KB)& 100 & 0.364 \\
\texttt{/createToDo} (file, 1\,MB)  & 100 & 0.732 \\
\texttt{/createToDo} (file, 5\,MB)  & 100 & 2.341 \\
\texttt{/login}                   & 100 & 6.750 \\
\midrule
\texttt{/ai} (very short prompt)  &  10 & 39.15  \\
\texttt{/ai} (short prompt)       &  10 & 313.97 \\
\texttt{/ai} (medium prompt)      &  10 & 373.86 \\
\texttt{/ai} (long prompt)        &  10 & 286.73 \\
\texttt{/ai} (very long prompt)   &  10 & 342.09 \\
\bottomrule
\end{tabular}
\end{table}

The results in Table~\ref{tab:energy} illustrate several distinct energy regimes.

\textbf{Constant-energy CRUD endpoints.} Simple write and update endpoints (\texttt{/done}, \texttt{/createToDo} with short text, \texttt{/logout}) cluster in the range 0.13--0.47\,mWh per request. These endpoints show little variation across conditions and are best modelled as constants. The system idle power measured at the PSU during the dedicated idle phase was $\approx$41\,W, confirming that most of this budget is attributable to baseline system draw rather than request processing.

\textbf{Payload-sensitive write endpoint.} For \texttt{/createToDo} with text payloads of 100 to 100\,000 characters, the per-request energy varies only modestly (0.38--0.47\,mWh), suggesting that for text-dominated payloads, network and serialisation costs are negligible at local scales. In contrast, file attachments show a clear non-linear increase: 1\,KB--100\,KB attachments remain flat near 0.37\,mWh, while 1\,MB jumps to 0.73\,mWh and 5\,MB reaches 2.34\,mWh. A piecewise model is appropriate for this endpoint: constant below roughly 100\,KB, linear above.

\textbf{State-dependent read endpoint.} \texttt{/getToDos} scales with the number of stored items returned: from 0.196\,mWh (100 items) to 1.300\,mWh (100\,000 items). The relationship is approximately linear up to 10\,000 items, with a steeper jump at 100\,000 items, consistent with a linear or curve model parameterised by result-set size.

\textbf{Authentication cost.} \texttt{/login} consumes 6.75\,mWh per request, an order of magnitude above comparable CRUD endpoints. This is attributable to the bcrypt password hashing used by the application, which is intentionally CPU-intensive. This illustrates a class of endpoints that are computationally heavy by design and should be modelled as constants with explicit documentation of their security-driven cost.

\textbf{LLM inference dominates.} AI inference requests are two to three orders of magnitude more expensive than any CRUD operation. Even a very-short prompt consumes 39\,mWh, while typical prompt lengths reach 287--374\,mWh per request. GPU power draw during the AI phases averaged 88.8\,W compared to an idle GPU draw of 17.6\,W, confirming that the GPU accounts for the bulk of the additional cost. Energy does not scale monotonically with prompt length (the long prompt is cheaper than the medium one), reflecting the non-deterministic token generation behaviour of the model. An approximate constant model captures the order of magnitude; a curve model parameterised by input or output token count would be required for finer-grained disclosure.

These results validate the three-tier model hierarchy introduced in Section~\ref{sec:method}: constant, linear, and curve models cover the full range of endpoint behaviours observed, and the energy disclosure values span five orders of magnitude across endpoint types---from sub-mWh CRUD calls to hundreds of mWh for LLM inference---underscoring that aggregating endpoints into a single server-level number would be misleading.


------


%-----------------------------------------------------------------------------------------
\section{Motivation}
%-----------------------------------------------------------------------------------------

\paragraph{Carbon accounting for software and HTTP request granularity.}
A central challenge in green software is to connect engineering observables with carbon-accounting concepts.  For HTTP-based systems, that functional unit is naturally the individual HTTP request, because this is the unit at which web services are invoked, counted, billed, rate-limited, and experienced by users. Applying SCI to HTTP services therefore implies a per-request perspective: the relevant question is not only how much carbon a service emits in total, but how much impact is associated with processing a single request. This shift is important because only request-level values can be surfaced at the HTTP boundary and used in practice for software decisions such as retries, batching, caching, provider selection, or prompt reduction. In this sense, SCI gives the accounting logic, and HTTP provides the interface at which that logic becomes operational. \cite{kern2015impacts,lago2015framing,verdecchia2021green}.

\paragraph{Why HTTP is a plausible disclosure surface?}
HTTP is an attractive layer for impact disclosure because it is already the dominant interoperability boundary for digital services. It is where costs are billed, quotas are enforced, latency is observed, and usage is counted. Exposing sustainability information at the same boundary would therefore make environmental impact legible at the point where technical and organisational decisions are made. More generally, HTTP explicitly supports extensibility through header fields \cite{rfc9110}, and Structured Field Values provide an established way to encode non-trivial metadata in interoperable form \cite{rfc8941}.

Environmental information only becomes operationally useful once it is surfaced at interfaces that software and operators can actually act upon. Treehouse argues for making energy and carbon visible to datacenter software developers at fine granularity \cite{anderson2023treehouse}, while \textit{carbond} proposes an operating-system daemon that mediates carbon information between hardware, power supply, and application-level software \cite{schmidt2024carbond}. These works are not HTTP-header proposals, but they support the theoretical premise that environmental information should be promoted to a first-class systems interface.

\paragraph{Previous work on HTTP carbon and sustainability headers.}
The most direct prior work on HTTP-level disclosure is Martin's 2023 Internet-Draft, which proposed a single response header (Carbon-Emissions-Scope-2) for reporting the emissions associated with processing a request and producing the corresponding response \cite{martin-http-carbon-emissions-scope-2-00}. The proposal is conceptually important because it establishes the basic idea that environmental impact can be surfaced directly in HTTP responses, analogously to other operational metadata. A more recent 2026 Internet-Draft expands this idea by proposing a structured \texttt{Sustainability} response header \cite{besleaga-green-sustainability-header-00}. In contrast to the single-value design of Martin's draft, this proposal uses a Structured Field representation and is explicitly designed to carry multiple sustainability dimensions, including Scope~2 and Scope~3 emissions.  What these proposals do \textbf{not} yet resolve is the server-side problem of how a provider should generate credible per-request values. A field name and syntax can standardise \emph{what} is communicated, but not \emph{how} the underlying quantity is obtained. For real systems, that is the harder problem.

\paragraph{Visibility and web-facing sustainability information.}
\textit{Green Web Meter} presents a real-time digital sustainability monitoring system for websites and web applications, showing that environmental information can be integrated into operational web tooling and exposed in ways that support assessment and improvement \cite{sala2024greenwebmeter}. CASPER shows that distributed web services can respond to carbon signals through carbon-aware scheduling and provisioning decisions while still respecting service-level constraints \cite{souza2023casper}.

\paragraph{API user (consumer) perspective}
For API consumers per request resource consumption matters because service usage is typically visible through price, latency, and quota consumption, but not through energy or \(\mathrm{CO_2e}\). When an application calls \texttt{/login}, \texttt{/search}, or \texttt{/ai}, the client can count requests and measure response times, yet it usually cannot observe the emissions associated with those calls. If impact were disclosed directly in the response, consumers could use it for carbon accounting, provider comparison, and software optimisation decisions such as batching, caching, retry behaviour, or prompt reduction. This is precisely the kind of per-functional-unit reasoning encouraged by SCI \cite{gsf-sci}.

\paragraph{Service provider perspective}
Providers face the complementary problem. They are increasingly expected to understand, reduce, and sometimes disclose the environmental impact of their services, but request-level attribution is difficult in contemporary deployments. Virtualisation, multi-tenancy, managed platforms, and opaque downstream dependencies all make direct power attribution difficult or impossible in production. At the same time, providers would benefit from a practical disclosure mechanism not only for external transparency but also for internal optimisation and observability. A disclosure mechanism at the HTTP layer could support both uses, provided that the values can be generated cheaply and consistently.

\paragraph{Why a model is needed}
Previous work on software sustainability measurement makes clear that the field still relies on a heterogeneous collection of setups, methods, and abstractions, and that a reference measurement model is needed to structure and compare approaches \cite{guldner2024gsmm}. For HTTP sustainability disclosure, this implies that direct online measurement is not a prerequisite for useful disclosure. Instead, what is needed is a defensible method that links controlled measurements to runtime-observable request features. The theoretical background leads to a clear gap. SCI provides the accounting logic, prior HTTP drafts provide candidate disclosure semantics, and related systems work motivates making carbon information visible at software interfaces. What remains missing is a practical bridge from \emph{measurement} to \emph{per-request disclosure}.


The central implementation challenge of per-request sustainability disclosure is the derivation of credible values. In principle, one might want to attach energy use or \(\mathrm{CO_2e}\) directly to every HTTP request as it is processed. In practice, however, these quantities are not directly observable at the request level. Power is usually measured only at the machine or system boundary, while HTTP requests are handled concurrently, interact with shared resources, and often trigger background work that is not trivially attributable to a single transaction.

This creates a gap between what the protocol intends to disclose and what the system can actually measure. A server can easily observe variables such as request duration, request bytes, response bytes, HTTP method, and response status, but it typically cannot directly observe the exact energy consumed by an individual request in production. Even where coarse power telemetry exists, it is affected by noise from idle power, caching, operating-system activity, virtualization, and interference from other workloads. As a result, raw power readings cannot simply be copied into HTTP headers.

A structured derivation process is therefore required. The role of the following steps is to bridge machine-level measurements and request-level disclosure in a transparent and reproducible way. The basic idea is to begin with controlled benchmark measurements, transform these measurements into per-request estimates, and then relate them to features that are available at runtime at the HTTP boundary. Only after this transformation is it possible to generate request-specific sustainability information online with low overhead. These runtime features are not limited to timing and byte counts: commodity web servers can also observe request method, response status, and selected properties of the request URI, including URL parameters, without requiring invasive instrumentation.

The steps below are necessary for two reasons. First, they make explicit how endpoint-level energy values are obtained from raw measurements instead of treating them as assumed constants. Second, they provide the methodological link between offline experimentation and online disclosure.
